\section{The canonical resolvent and its exact kernel series}
\label{canonical:section}

Theorem~\ref{erc:equality} gives $W(s)<\infty$ for every $s>\mu$.
Pass to the limit $T\to\infty$ in Theorem~\ref{erc:gradient} to obtain
\begin{equation}\label{erc:infinite-gradient}
 |V_s(y^+)-V_s(y)|\le\frac{40}{s}W(s),\qquad s>\mu.
\end{equation}
A certified block $e^{-sT}H(T)\le q<1$ supplies the explicit bound
\begin{equation}\label{canonical:W-bound}
 W(s)\le\int_0^\infty e^{-st}H(t)\,dt
 \le\frac1{1-q}\int_0^T e^{(1-s)t}\,dt.
\end{equation}
Such a block can be found by the complete search after
Theorem~\ref{sc:window}, now with $\lambda=\mu$.

The old-core decomposition gives
\begin{equation}\label{canonical:linear}
 \frac{C(y)}s\le V_s(y)
 \le\frac{C(y)}s(1+W(s)).
\end{equation}
For the upper inequality, partition the future population by the
currently existing cores and their future newborn families, whose
birth intensities are at most one. Integrating their discounted
rewards gives the right side. The lower inequality uses monotonicity
of the core count. These estimates, finite-horizon population moments,
and localization justify the resolvent generator equation
\begin{equation}\label{canonical:generator}
 (s-\mathcal L_\infty)V_s=C.
\end{equation}

The block also controls the remaining time integral effectively:
\begin{equation}\label{canonical:time-tail}
 \int_{nT}^\infty e^{-st}P_t^\infty C(y)\,dt
 \le\frac{C(y)q^n}{1-q}\int_0^T e^{(1-s)r}\,dr.
\end{equation}
Together with a finite-state calculation on $[0,nT]$, this computes
$V_s(y)$ at every finite state, once a block with $e^{-sT}H(T)<1$ has been certified.

\begin{proposition}[Exact eliminated-kernel representation]
\label{canonical:series}
For $s>\mu$, let $Q_s$ be the exact core-ring kernel in
\eqref{model:Q}, and let
\[
 g_s(y)=\frac{C(y)}{C(y)+s}.
\]
Then
\begin{equation}\label{canonical:kernel-series}
 V_s(y)=\sum_{j=0}^\infty(Q_s^jg_s)(y).
\end{equation}
This is a nonnegative convergent series with explicitly given beta
kernel entries. Its value is effectively computable at every finite
state and rational $s>\mu$. In particular $V_s$ is the canonical
positive supersolution for the eliminated kernel, with the uniform
root-leaf increment \eqref{erc:infinite-gradient}.
\end{proposition}
\begin{proof}
Let $\tau_j$ be the successive core-ring times, with $\tau_0=0$.
Between core rings $C$ is constant. The expected discounted reward
before the first ring is $R_sC=g_s$, and the discounted law of the
post-ring state is $Q_s$. The Markov property identifies the expected
reward over $[\tau_j,\tau_{j+1})$ with $Q_s^jg_s$. The original
root's rate-one clock ensures infinitely many rings over an infinite
time interval; nonexplosion ensures $\tau_j\to\infty$. Tonelli's
theorem partitions the total reward into these intervals and proves
\eqref{canonical:kernel-series}. Finiteness follows from
\eqref{canonical:linear}; effective evaluation follows from
\eqref{canonical:time-tail} and finite-state uniformization. The
first-step equation is $V_s=g_s+Q_sV_s$, and
\eqref{erc:infinite-gradient} gives the boundary increment.
\end{proof}

The finite-window construction also gives a computable supersolution
with bounded increments at \emph{every} ordinary birth, rather than
only at a root-leaf addition. At every rational $s>\mu$, a complete
search finds $T,q$ with $e^{-sT}H(T)\le q<1$, and
Theorem~\ref{sc:window} gives $U_{s,T}$ and its explicit increment
constant. The canonical and finite-window constructions are both
useful: the former solves the exact resolvent equation, while the
latter has a direct finite-time formula and controls all ordinary
birth increments.
